\documentclass[10pt]{article} 
\renewcommand{\thempfootnote}{\arabic{mpfootnote}}

\usepackage[utf8]{inputenc}

\usepackage{multicol}
\setlength{\columnsep}{1cm}

\usepackage{abc}
\usepackage{graphicx}
\usepackage{musicography}

\usepackage[margin=0.5in]{geometry} 
\geometry{letterpaper}

\usepackage{fancyhdr} 
\pagestyle{fancy} % options: empty , plain , fancy
\renewcommand{\headrulewidth}{0pt} % customise the layout...
\lhead{}\chead{}\rhead{}
\lfoot{}\cfoot{\thepage}\rfoot{}

\begin{document}

\thispagestyle{empty}

\begin{abc}[name=sevenup,options={-O= +c -F single.fmt}]
X:1
T:Seven Up
T:(Old Sledge)
S:Greg & Jere Canote, Portland jam session
Z:abc-transcription Josh Larios <hades@elsewhere.org>, 2021.05.29
Z:abc-copyright This transcription is licensed under a Creative Commons Attribution-ShareAlike 4.0 International License.
M:C|
K:G
|: Bc | "G"d2 d2 "C"e2-ee | "G"d2 d2 BA G2 | "D"A2 G2 AB c2 | "G"B2 G2 "D"A2 D2 | 
"G"d2 d2 "C"e2-ee | "G"d2 d2 BA G2 | "D"A2 G2 ABcA | BGAF "G"G2-GG | G2 :|
|: (3GFE | "G"D2 D2 "C"E2-E E | "G"D2 D2 G4 | D2 D2 "C"E2 ED | "G"B,A,G,B, "D"A,4 | 
"G"D2 D2 "C"E2-E E | "G"D2 D2 G2-GA | "D"B2 G2 ABcA | BGAF "G"G2-GG | G2 :|
\end{abc}
\vskip 0.5in
\noindent\makebox[\textwidth][c]{%
  \begin{minipage}{7.0in}
   \setlength{\parindent}{1.5em}

This tune comes from the Canote Brothers, who credit it to Ralph
Roberts of West Virginia, by way of Jimmy Triplett. Josh learned it in
a jam session at the Portland Old Time Music Gathering a few years
ago, and it's on the Canotes' CD with Julia Weatherford, ``Down in
North Carolina''.

``Seven Up'' and ``Old Sledge'' are alternate names for the traditional
English card game ``All Fours'', a trick-taking game first documented in
1674 and popular in North America through the 19th century. Ralph
Roberts apparently knew the tune primarily as ``Seven Up'', but also
called it ``Old Sledge''. This is a completely different tune than the
one more commonly known as ``Old Sledge'' played out of AEAD
tuning. It's also unrelated to the tunes named ``Seven Up'' found in
Ryan's Mammoth Collection and Marion Thede's ``The Fiddle Book'' (which
are unrelated to each other).

Perhaps it would be best to start calling this one ``All Fours'' to
distinguish it from the others, although that would be an ironic title
for a crooked tune. As it stands, though, it's nearly impossible to
find recordings of this particular tune, thanks to the overlapping
titles.

This is a stripped down transcription; listening to a recording by the
Canote Brothers is highly recommended. If you can't find a copy of
``Down in North Carolina'', there is a recording of them playing it for
their string band class at http://stringband.mossyroof.com/

   
 \end{minipage}
}


\end{document}
